\section{Betrieb}

\subsection{Start der Anwendung}

Für den Entwicklungsbetrieb ist \textbf{Docker Compose} der primäre
Startweg, da Frontend, Backend und Persistenzpfad in einem reproduzierbaren
Setup zusammengeführt werden.

\paragraph{Betrieb mit Docker Compose}
Für die Abgabe wird der Betrieb über die bereitgestellten Docker-Skripte
empfohlen, da diese die Konfiguration und den Start reproduzierbar durchführen.

\begin{lstlisting}[style=terminal,caption={Erstinitialisierung (Linux/macOS)}]
./docker_setup.sh
\end{lstlisting}

\begin{lstlisting}[style=terminal,caption={Regulaerer Start/Stop (Linux/macOS)}]
./docker_start.sh
./docker_shutdown.sh
\end{lstlisting}

Unter Windows stehen funktional gleichwertige Batch-Dateien bereit:

\begin{lstlisting}[style=terminal,caption={Windows (CMD/PowerShell)}]
docker_setup.bat
docker_start.bat
docker_shutdown.bat
\end{lstlisting}

Alternativ ist der manuelle Compose-Start weiterhin möglich:

\begin{lstlisting}[style=terminal]
cp .env.example .env
docker compose up
\end{lstlisting}

Das Anlegen der \texttt{.env} ist verpflichtend, da das Backend den Betrieb mit
einem unsicheren oder fehlenden \texttt{JWT\_SECRET} gezielt verweigert.
Im Setup-Skript kann \texttt{JWT\_SECRET} leer gelassen werden; in diesem Fall
wird automatisiert ein zufälliger, sicherer Secret-Wert generiert und in die
\texttt{.env} übernommen.

Nach dem Start stehen die zentralen Endpunkte bereit:
\begin{itemize}
  \item Frontend: \url{http://localhost:5173}
  \item Backend: \url{http://localhost:3000}
  \item Health-Endpoint: \url{http://localhost:3000/api/health}
\end{itemize}

Im Compose-Betrieb führt der Backend-Service vor dem Start automatisch
\texttt{prisma migrate deploy} aus.
Damit ist sichergestellt, dass das Datenbankschema dem aktuellen
Migrationsstand entspricht.
Das Setup-Skript ergänzt dies um die initiale Installation der Abhängigkeiten in
den Docker-Volumes (\texttt{npm ci} für Frontend und Backend), wodurch ein
frischer Pull ohne manuelle Zwischenschritte lauffähig wird.

\paragraph{Lokaler Betrieb ohne Docker}
Alternativ kann die Anwendung lokal in zwei Prozessen betrieben werden:

\begin{lstlisting}[style=terminal,caption={Backend lokal starten}]
cd backend
npm ci
npx prisma migrate deploy
npm run dev
\end{lstlisting}

\begin{lstlisting}[style=terminal,caption={Frontend lokal starten (zweites Terminal)}]
cd frontend
npm ci
npm run dev
\end{lstlisting}

Im lokalen Betrieb ist ein \texttt{backend/.env} mit den notwendigen Variablen
erforderlich (siehe \cref{subsec:betrieb-konfiguration}).

\subsection{Konfiguration}
\label{subsec:betrieb-konfiguration}

Die Anwendung wird über Umgebungsvariablen konfiguriert.
Für den Compose-Betrieb dient \texttt{.env.example} als Vorlage.

\begin{table}[h]
\centering
\small
\setlength{\tabcolsep}{4pt}
\renewcommand{\arraystretch}{1.3}
\begin{tabularx}{\textwidth}{>{\raggedright\arraybackslash}p{3.1cm} >{\raggedright\arraybackslash}p{3.4cm} >{\raggedright\arraybackslash}X}
\toprule
\textbf{Variable} & \textbf{Beispielwert} & \textbf{Bedeutung} \\
\midrule
\texttt{PORT} & \texttt{3000} & Port, auf dem das Backend lauscht (alternativ über Host vorgegeben). \\
\texttt{JWT\_SECRET} & \texttt{replace-with-a-long-random-secret} & Secret für Signierung und Verifikation der JWTs im Backend (Pflichtwert, unsichere Defaults sind blockiert). \\
\texttt{CORS\_ORIGIN} & \texttt{http://localhost:5173} & Erlaubte Origin(s) für CORS-Requests auf die API. \\
\texttt{DATABASE\_URL} & \texttt{file:./data/dev.db} & SQLite-Datenbankpfad des Backends. \\
\shortstack[l]{\texttt{VITE\_API\_}\\\texttt{PROXY\_TARGET}} & \texttt{http://backend:3000} & Ziel für den \texttt{/api}-Proxy des Vite-Dev-Servers. \\
\bottomrule
\end{tabularx}
\caption{Relevante Umgebungsvariablen im Entwicklungsbetrieb}
\label{tab:betrieb-env}
\end{table}

\paragraph{Voraussetzungen}
Für den Betrieb werden Node.js und npm benötigt.
Im Docker-Modus sind zusätzlich Docker Engine und Docker Compose erforderlich.
Die Persistenz erfolgt auf SQLite-Basis und benötigt keinen separaten
Datenbankserver.
Für die Windows-Skripte wird zusätzlich eine Standardinstallation von
PowerShell vorausgesetzt (zur Erzeugung eines zufälligen JWT-Secrets).

\subsection{Betriebsmodus}

Die Anwendung ist als \textbf{Entwicklungs- und Demo-Betrieb} ausgelegt.
Frontend und Backend laufen als getrennte Prozesse/Container, kommunizieren
über HTTP und nutzen ein Cookie-basiertes Session-Modell mit JWT.

\paragraph{Typischer Workflow}
\begin{enumerate}
  \item Services starten (Compose oder lokal).
  \item Benutzer über \texttt{/register} anlegen und anmelden.
  \item Campaign erstellen oder per Join-Code beitreten.
  \item Inhalte pflegen (abhängig von der Rolle DM/EDITOR/PLAYER).
\end{enumerate}

\paragraph{Betriebsrelevante Eigenschaften}
\begin{itemize}
  \item Daten werden persistent in SQLite gespeichert.
  \item Schemaänderungen werden über Prisma-Migrationen ausgerollt.
  \item Geschützte Endpunkte erfordern gültige JWT-Cookies.
  \item Aktive Sitzungen werden über \texttt{/api/session/refresh} kontrolliert
  verlängert.
  \item Der Health-Endpoint \texttt{/api/health} liefert zusätzlich \texttt{version} und \texttt{uptimeSeconds}.
  \item Die Trennung von Frontend und Backend erlaubt klare Fehlersuche
  (UI-/API-/DB-Ebene).
\end{itemize}
